\begin{summarybox}
\textbf{一条主线：焦角 $\longleftrightarrow$ 面积 $\longleftrightarrow$ 高度
$\longrightarrow$ 坐标绝对值。}

\begin{enumerate}
  \item \textbf{核心通式（任意点）：}
  \[
  \boxed{\tan\frac\theta2=\frac{c|y_0|}{b^2}},
  \qquad 0\leq|y_0|\leq b.
  \]

  \item \textbf{最大角（短轴端点）：}
  \[
  \boxed{P=(0,\pm b)},\qquad
  \boxed{\tan\frac{\theta_{\max}}2=\frac cb},\qquad
  \boxed{\theta_{\max}=2\arcsin e}.
  \]

  \item \textbf{最小角（长轴端点）：}
  \[
  \boxed{P=(\pm a,0)},\qquad
  \boxed{\theta_{\min}=0^\circ},\qquad
  \boxed{0^\circ\leq\theta\leq\theta_{\max}<180^\circ}.
  \]

  \item \textbf{面积、焦角、高度三者等价（固定椭圆）：}
  \[
  \boxed{S=c|y_0|=b^2\tan\frac\theta2},\qquad
  \boxed{|x_0|=a\sqrt{1-\frac{y_0^2}{b^2}}}.
  \]
  焦角确定，两个坐标的绝对值确定；\textbf{一般有四个对称点，
  最大角与最小角各只有两个点}。

  \item \textbf{余弦定理双重核验：}
  \[
  \boxed{\cos\theta
   =\frac{r_1^2+r_2^2-4c^2}{2r_1r_2}
   =\frac{b^4-c^2y_0^2}{b^4+c^2y_0^2}}.
  \]

  \item \textbf{两种记法：}通式记成
  ``$c\times\text{高}\div b^2$''；最大值记成 ``$c\div b$''。
  \textbf{千万不要把 $c/|y_0|$ 记成任意点公式。}

  \item \textbf{证明思路：}余弦定理 $\to$ 半角公式；
  或面积 $S=c|y_0|$ $\to$ 内切圆半径 $\to$ 半角正切；
  还可由 $r_1+r_2=2a$ 和余弦定理先得
  $r_1r_2=b^2/\cos^2(\theta/2)$，再直接求面积
  $S=b^2\tan(\theta/2)$。
  直角三角形 $\tan(\theta_{\max}/2)=c/b$ 只在短轴端点直接使用。
\end{enumerate}
\end{summarybox}
